\section{Techniques}

\subsection{Technique 1: Purely Functional Data Structure}

\textbf{Module:} \texttt{src/ScopeStack.hs}

The central data structure is the \emph{scope stack}: a persistent,
purely functional stack of hash-map frames.

\begin{figure}[ht!]
    \centering
    \begin{minipage}{.9\linewidth}
\begin{lstlisting}[language=Haskell]
newtype ScopeStack a = ScopeStack [Map.Map Ident a]

push :: ScopeStack a -> ScopeStack a
push (ScopeStack fs) = ScopeStack (Map.empty : fs)

pop :: ScopeStack a -> ScopeStack a
pop (ScopeStack (_:fs@(_:_))) = ScopeStack fs
pop (ScopeStack _)             = ScopeStack [Map.empty]

lookupStack :: Ident -> ScopeStack a -> Maybe a
lookupStack x (ScopeStack fs) = foldr (\f acc ->
  case Map.lookup x f of { Just v -> Just v; _ -> acc }) Nothing fs

updateStack :: Ident -> a -> ScopeStack a -> ScopeStack a
updateStack x v (ScopeStack fs) = ScopeStack (go fs)
  where go []     = []
        go (f:rest)
          | Map.member x f = Map.insert x v f : rest
          | otherwise      = f : go rest
\end{lstlisting}
    \end{minipage}
    \captionof{lstlisting}{\texttt{ScopeStack}: push, pop, lookup, update}
    \label{lst:scopestack}
\end{figure}

\noindent
\textbf{Why it qualifies as a persistent data structure.}
\texttt{push} and \texttt{pop} are O(1) list cons and tail --- the old stack
is never overwritten.
After a \texttt{pop}, the caller still holds a reference to the original
\texttt{push}ed version; no copy is needed.
\texttt{updateStack} rewrites only the \emph{first frame} that contains the
key and shares all other frames with the original, so two callers holding
different versions of the stack remain consistent.

\noindent
\textbf{Why \texttt{updateStack} matters for ownership.}
When an inner block assigns to an outer variable
(\texttt{let mut x = 0; \{ x = 5 \}}),
\texttt{updateStack} finds \texttt{x} in an outer frame and updates it
\emph{in place within that frame}, while the inner scope frame is discarded
on block exit.
Without this operation the persistent structure would silently ignore the
mutation.

\noindent
\textbf{Advanced usage.}
The same \texttt{ScopeStack} is the \emph{sole mechanism} for variable
scoping, mutability, ownership tracking, borrow counting, and lifetime
management in every phase.
Both the type checker (\texttt{TcCtx.\_tcVars :: ScopeStack VarInfo}) and
the interpreter (\texttt{EvalCtx.\_evalVars :: ScopeStack Value}) use it.
There is no separate ownership map, no borrow graph, and no auxiliary data
structure for scoping.

% ---------------------------------------------------------------------------
\subsection{Technique 2: Lenses and Traversals}
% ---------------------------------------------------------------------------

\textbf{Modules:} \texttt{src/Context.hs} (lenses),
                  \texttt{src/TypeCheck/Stmt.hs} (NLL traversal)

\subsubsection*{Van Laarhoven Lenses}

The project uses five hand-rolled van Laarhoven lenses to focus into the
nested fields of \texttt{TcCtx} and \texttt{EvalCtx}:

\begin{figure}[ht!]
    \centering
    \begin{minipage}{.9\linewidth}
\begin{lstlisting}[language=Haskell]
type Lens' s a = forall f. Functor f => (a -> f a) -> s -> f s

view :: Lens' s a -> s -> a
view l s = getConst (l Const s)

over :: Lens' s a -> (a -> a) -> s -> s
over l f s = runIdentity (l (Identity . f) s)

set :: Lens' s a -> a -> s -> s
set l v = over l (const v)

-- Example lens for the variable scope in the type-checker context:
tcVars :: Lens' TcCtx (ScopeStack VarInfo)
tcVars f ctx = (\v -> ctx { _tcVars = v }) <$> f (_tcVars ctx)
\end{lstlisting}
    \end{minipage}
    \captionof{lstlisting}{Hand-rolled van Laarhoven lens type and combinators}
    \label{lst:lenses}
\end{figure}

\noindent
Every TypeCheck and Interp module uses \texttt{modify (over tcVars ...)} to
update nested context fields.
Without lenses, each update would require a manual record-update expression
and repeated \texttt{gets} calls.

\subsubsection*{Traversal as a Structural Fold (NLL)}

The advanced usage is in \texttt{releaseExpiredBorrows} in
\texttt{TypeCheck/Stmt.hs}.
This function must identify borrows whose holder is no longer live after a
statement, but it must not modify the scope stack --- it only needs to
\emph{collect} the expired entries.

\begin{figure}[ht!]
    \centering
    \begin{minipage}{.9\linewidth}
\begin{lstlisting}[language=Haskell]
releaseExpiredBorrows :: Set.Set Ident -> Tc ()
releaseExpiredBorrows live = do
  vars <- gets (view tcVars)
  let expired = getConst $
        Map.traverseWithKey
          (\x vi -> Const $
            if x `Set.notMember` live && isJust (varBorrowOf vi)
            then [(x, vi)]
            else [])
          (topBindings vars)
  mapM_ releaseAndClear expired
\end{lstlisting}
    \end{minipage}
    \captionof{lstlisting}{NLL fold using \texttt{Const} applicative as a structural fold}
    \label{lst:nll-fold}
\end{figure}

\noindent
\textbf{Why this is a traversal.}
\texttt{Map.traverseWithKey} has type
\texttt{Applicative f => (k -> a -> f a) -> Map k a -> f (Map k a)}.
It is a van Laarhoven traversal over the \texttt{(key, value)} pairs of a map.
When the applicative \texttt{f} is instantiated to
\texttt{Const [(Ident, VarInfo)]}, the \texttt{<*>} operation is list
concatenation (from the \texttt{Monoid} instance of lists), and no new map is
ever allocated --- the ``return value'' is discarded via \texttt{getConst}.
This transforms a structure-mapping operation into a
\emph{zero-allocation structural fold}, collecting all expired borrow entries
in a single traversal of the top scope frame.
The same function instantiated with \texttt{Identity} would \emph{update}
the map; with \texttt{Const} it only \emph{reads} from it.

% ---------------------------------------------------------------------------
\subsection{Technique 3: Monad Transformers}
% ---------------------------------------------------------------------------

\textbf{Modules:} \texttt{src/Tc.hs}, \texttt{src/Eval.hs}

\begin{figure}[ht!]
    \centering
    \begin{minipage}{.9\linewidth}
\begin{lstlisting}[language=Haskell]
-- Type-checker monad
type Tc a = ExceptT String (State TcCtx) a

runTc :: TcCtx -> Tc a -> Either String a
runTc ctx m = evalState (runExceptT m) ctx

-- Interpreter monad
type Eval a = ExceptT String (State EvalCtx) a

runEval :: EvalCtx -> Eval a -> Either String a
runEval ctx m = evalState (runExceptT m) ctx
\end{lstlisting}
    \end{minipage}
    \captionof{lstlisting}{\texttt{Tc} and \texttt{Eval} monad transformer stacks}
    \label{lst:monads}
\end{figure}

\noindent
\textbf{Why \texttt{ExceptT} over \texttt{State}.}
The type checker needs two capabilities simultaneously:
\begin{itemize}
  \item \emph{Mutable state}: the context $\Gamma$ is updated as bindings are
    introduced, variables are moved, and borrow counts change.
    \texttt{State TcCtx} provides \texttt{gets} and \texttt{modify}.
  \item \emph{Short-circuiting errors}: a type error should abort the current
    check immediately without running further checks.
    \texttt{ExceptT String} provides \texttt{throwError}.
\end{itemize}

Neither capability alone suffices.
A pure \texttt{Either String (TcCtx, a)} function would need to thread the
context manually through every call.
A \texttt{State TcCtx (Either String a)} would not short-circuit --- errors
would be accumulated as \texttt{Left} values without stopping execution.
The transformer stack \texttt{ExceptT String (State TcCtx)} gives both:
\texttt{throwError} short-circuits and the state persists until then.

\begin{figure}[ht!]
    \centering
    \begin{minipage}{.9\linewidth}
\begin{lstlisting}[language=Haskell]
-- Typical usage in TypeCheck/Stmt.hs
infer (SAssign x e) = do
  vars <- gets (view tcVars)             -- read state
  case lookupStack x vars of
    Nothing -> throwError "not in scope" -- short-circuit on error
    Just vi -> do
      unless (varMut vi) $
        throwError "Cannot assign to immutable variable"
      E.check e (varType vi)
      modify (over tcVars (updateStack x vi { varOwned = True }))
      -- persist state update ^
\end{lstlisting}
    \end{minipage}
    \captionof{lstlisting}{Typical use of \texttt{gets}, \texttt{throwError}, and \texttt{modify}}
    \label{lst:monad-usage}
\end{figure}

\noindent
The two stacks (\texttt{Tc} and \texttt{Eval}) are structurally identical;
only the context type differs (\texttt{TcCtx} vs.\ \texttt{EvalCtx}).
This means every type-checker combinator (\texttt{throwError}, \texttt{gets},
\texttt{modify}) has an exact counterpart in the interpreter, making both
codebases read identically at the monad-plumbing level.

% ---------------------------------------------------------------------------
\subsection{Technique 4: Concurrency}
% ---------------------------------------------------------------------------

\textbf{Modules:} \texttt{src/Logger.hs} (infrastructure),
\texttt{src/TypeCheck/Stmt.hs} + \texttt{src/Interp/Stmt.hs} (language-level)

Concurrency is used at two levels: an \emph{infrastructure} level (a background
logging thread) and a \emph{language} level (a user-facing \texttt{spawn}
construct whose safety is enforced by the type checker).

\subsubsection*{Background Logger (\texttt{forkIO} + \texttt{Chan} + \texttt{MVar})}

\begin{figure}[ht!]
    \centering
    \begin{minipage}{.9\linewidth}
\begin{lstlisting}[language=Haskell]
data Logger = Logger
  { _chan    :: Chan (Maybe LogMsg)  -- message channel; Nothing = stop
  , _done    :: MVar ()             -- drain thread signals completion
  , _enabled :: Bool
  }

startLogger :: Bool -> IO Logger
startLogger enabled = do
  ch   <- newChan
  done <- newEmptyMVar
  _    <- forkIO (drain ch done)    -- fork background thread
  return (Logger ch done enabled)

stopLogger :: Logger -> IO ()
stopLogger (Logger ch done _) = do
  writeChan ch Nothing              -- send stop sentinel
  takeMVar done                     -- block until drain completes
\end{lstlisting}
    \end{minipage}
    \captionof{lstlisting}{Logger using \texttt{forkIO}, \texttt{Chan}, and \texttt{MVar}}
    \label{lst:logger}
\end{figure}

\noindent
The background thread drains the \texttt{Chan} until it receives a
\texttt{Nothing} sentinel, then signals completion via the \texttt{MVar}.
This is a classic producer-consumer pattern: the main thread writes messages
asynchronously; the drain thread prints them; \texttt{stopLogger} provides
a synchronous barrier that prevents messages from being lost on program exit.
Five logger tests verify that the presence of the background thread does not
affect type-checker or interpreter results.

\subsubsection*{Language-Level \texttt{spawn}}

\begin{figure}[ht!]
    \centering
    \begin{minipage}{.9\linewidth}
\begin{lstlisting}[language=Haskell]
-- Type checker: enforce Copy-only captures
infer (SSpawn body) = do
  let freeVs = mentionedBlock body
  vars <- gets (view tcVars)
  let checkCapture x = case lookupStack x vars of
        Nothing -> return ()
        Just vi -> unless (isCopyable (varType vi)) $ throwError $
          "Cannot capture non-Copy variable '" ++ printTree x
          ++ "' in spawn block"
  mapM_ checkCapture (Set.toList freeVs)
  checkBlock body

-- Interpreter: run synchronously (safety guarantee is in the type checker)
interp (SSpawn body) = interpBlock body
\end{lstlisting}
    \end{minipage}
    \captionof{lstlisting}{\texttt{spawn} type checking and interpretation}
    \label{lst:spawn}
\end{figure}

\noindent
\textbf{Thread-safety via the type system.}
The \texttt{isCopyable} predicate classifies \texttt{int}, \texttt{bool}, and
immutable references (\texttt{\&T}) as \texttt{Copy}; \texttt{Light}, lists,
\texttt{Result} of non-\texttt{Copy}, and mutable references
(\texttt{\&mut T}) are not \texttt{Copy}.
Because \texttt{Copy} types have value semantics (no aliasing), sharing them
between threads is always safe.
Non-\texttt{Copy} types have aliasing that would constitute a data race, so
they are rejected at the \texttt{spawn} site.
The type checker thus statically guarantees that a real concurrent execution
would be data-race-free --- no runtime locks are needed.

\noindent
\textbf{Two-level story.}
The background logger demonstrates low-level concurrent infrastructure
(fork, channel, synchronisation barrier).
The language-level \texttt{spawn} demonstrates that the \emph{ownership type
system} itself is the thread-safety mechanism: the affine and copy disciplines
that prevent use-after-move also prevent data races, without any additional
annotations.
